\chapter{Git and GitHub: evidence, collaboration, and recovery}
\label{app:git-github}

Git is sometimes introduced as a place to save code and GitHub as a place to
upload it. Both descriptions are too weak. A professional data project changes
continuously: code is revised, datasets acquire new versions, notebooks are
rerun, assumptions are corrected, and several people may work at once. Without
an inspectable history, the team cannot reliably answer which change produced a
result, why a decision was made, whether two analyses used the same source, or
how to return to the last known good state.

Git makes coherent project states identifiable and comparable. GitHub adds a
shared location and collaboration system around that history. Together they
support review, attribution, automation, release, and recovery. They do not
replace backups, data governance, testing, or scientific records; they connect
those practices to concrete versions of the project.

\begin{learningobjectives}
After working through this appendix, you should be able to:
\begin{itemize}
  \item distinguish Git, GitHub, a repository, a commit, a branch, a remote,
  and a pull request;
  \item explain the working tree, staging area, local repository, and remote
  repository without relying on memorized commands;
  \item inspect, stage, commit, compare, and recover changes deliberately;
  \item use a branch-based GitHub workflow for a small reviewed change;
  \item synchronize with collaborators and resolve a simple merge conflict;
  \item protect secrets, large data, environments, and generated artifacts;
  \item interpret a GitHub Actions workflow and use CI evidence during review;
  and
  \item follow a conservative daily workflow that works in VS Code on Windows,
  macOS, and Linux.
\end{itemize}
\end{learningobjectives}

\begin{chapterconnection}
Version control is part of the professional-practice spine, not an isolated
tool. It identifies the code and documentation used by notebooks, packages,
tests, data pipelines, models, APIs, and deployments. Chapter~5 develops tests
and CI; this appendix supplies the history and review workflow that make those
checks useful during change.
\end{chapterconnection}

\section{Why version control matters in data science}

A data-science result is rarely produced by one immutable script. Exploration
creates alternatives; cleaning rules evolve; features are redefined; model
parameters change; plots are revised; and conclusions are challenged. The
history matters for at least six reasons.

\begin{description}
  \item[Reproducibility.] A result can name the exact source revision that
  produced it rather than ``the code from last month.''
  \item[Review.] A collaborator can inspect the proposed difference instead of
  rereading the entire project and guessing what changed.
  \item[Attribution.] Commits record authorship, time, and intent. They help a
  team locate the person or discussion with relevant context; they are not a
  measure of individual productivity.
  \item[Parallel work.] Branches allow independent changes to proceed without
  continuously overwriting one shared directory.
  \item[Recovery.] Recorded snapshots provide known states to compare, restore,
  or counteract when an experiment or release fails.
  \item[Automation.] CI systems can test the precise revision proposed for
  integration and attach the evidence to its review.
\end{description}

Version control is evidence about artifacts and decisions. It does not prove
that the data was valid, the analysis was unbiased, or the conclusion was true.
A perfectly versioned invalid experiment remains invalid. The value of Git is
that the exact artifact can be examined, challenged, and connected to other
evidence.

\begin{misconceptionbox}
\textbf{Misconception:} Git is useful only when several programmers collaborate.

\textbf{Correction:} a solo researcher also changes assumptions, revisits old
results, makes mistakes, and needs to communicate with a future version of
themself. Small coherent commits make that reasoning inspectable before a second
person joins the project.
\end{misconceptionbox}

\section{Git and GitHub are different systems}

\term{Git} is an open-source distributed version-control system that runs on
your computer. Most Git operations, including commits, branches, comparisons,
and history inspection, do not require a network connection. A clone contains a
project history and can create new history locally.

\term{GitHub} is a hosted collaboration platform built around Git repositories.
It provides remote hosting, accounts and permissions, pull requests, reviews,
issues, releases, Actions, and other services. GitHub is not required to use
Git, and Git can communicate with hosting services other than GitHub.

\begin{center}
\small
\begin{tabularx}{\textwidth}{@{}
  >{\raggedright\arraybackslash}p{0.22\textwidth}
  >{\raggedright\arraybackslash}X
  >{\raggedright\arraybackslash}X@{}}
\toprule
\textbf{Question} & \textbf{Git} & \textbf{GitHub} \\
\midrule
Where does it operate? & Primarily in a local repository on your computer & As
a network service containing hosted repositories and collaboration records \\
What does it identify? & Snapshots, commits, branches, tags, and relationships
between them & Accounts, repositories, pull requests, reviews, issues, workflow
runs, releases, and permissions \\
Does it require the internet? & Not for ordinary local history operations & Yes
for interacting with the hosted service \\
Typical commands or actions & \code{status}, \code{diff}, \code{add},
\code{commit}, \code{switch}, \code{merge}, \code{log} & Open a pull request,
review files, inspect checks, manage an issue, merge an approved change \\
What does ``save'' mean? & Writing a file is separate from committing a staged
snapshot & Pushing transfers local commits; the GitHub copy is not updated by a
local save or commit alone \\
\bottomrule
\end{tabularx}
\end{center}

VS Code supplies a graphical Source Control view for many Git operations, but
the underlying states are the same. This appendix shows commands because they
make inputs and effects explicit and transfer across operating systems. You may
use the graphical interface after you can predict which Git state it will
change.

\begin{workedexample}{Saved is not committed, and committed is not shared}
A student edits \code{analysis.py} in VS Code and saves it. The file is now on
disk in the working tree, but no Git history contains the edit. The student runs
\code{git add analysis.py}; the current contents enter the staging area, but no
commit exists yet. After \code{git commit}, the snapshot exists in the local
repository. A collaborator still cannot see it on GitHub until the containing
branch is pushed to a remote the collaborator can access.

Each claim has distinct evidence: the editor or filesystem shows the saved file;
\code{git diff} and \code{git status} show the working change;
\code{git diff --staged} shows the proposed snapshot; \code{git log} shows the
local commit; and the GitHub branch or a fetch from the remote shows that the
commit was shared. Saying ``I saved it'' leaves all but the first state unknown.
\end{workedexample}

\section{The four-state mental model}

Most beginner confusion disappears when four locations are kept distinct.

\begin{center}
\begin{tikzpicture}[node distance=10mm]
  \node[flowbox] (work) {Working tree\\files you edit};
  \node[flowbox,right=of work] (stage) {Staging area\\next snapshot};
  \node[flowbox,right=of stage] (local) {Local repository\\commits};
  \node[flowbox,right=of local] (remote) {Remote repository\\shared commits};
  \draw[flowarrow] (work) -- (stage);
  \draw[flowarrow] (stage) -- (local);
  \draw[flowarrow] (local) -- (remote);
  \draw[flowarrow] (remote.south) .. controls +(0,-0.65) and +(0,-0.65) .. (local.south);
\end{tikzpicture}

\smallskip
{\small
\code{git add}: working tree $\rightarrow$ staging \quad
\code{git commit}: staging $\rightarrow$ local repository\\[2pt]
\code{git push}: local $\rightarrow$ remote \quad
\code{git fetch}: remote $\rightarrow$ local
}
\end{center}

\textbf{The working tree} is the checked-out directory visible in VS Code. Saving
a file changes the working tree, not Git history.

\textbf{The staging area}, also called the index, is Git's proposed content for
the next commit. Staging is selection, not a permanent save. If a staged file is
edited again, one version can be staged while newer edits remain unstaged.

\textbf{The local repository} is the history stored under the hidden
\code{.git} directory. A commit records the staged project snapshot, metadata,
a message, and parent commit. It does not automatically appear on GitHub.

\textbf{A remote repository} is another repository reachable through a named
network location. \code{origin} is a conventional local nickname for the remote
used when cloning; it is not a special server or a branch.

These locations explain common surprises. \code{git diff} normally shows
unstaged working-tree changes. \code{git diff --staged} shows what the next
commit would contain. A clean working tree says that tracked local files match
the current commit; it does not say the branch has been pushed or is current
with GitHub.

\begin{professionalbox}
Use \code{git status} before and after every unfamiliar Git command. Read its
output as evidence: current branch, relationship to its upstream branch, staged
changes, unstaged changes, and untracked files. Do not treat status as an error
message to dismiss.
\end{professionalbox}

\section{Commits, branches, and HEAD}

Git stores project snapshots as commits connected to parent commits. The
resulting history is a directed graph. A branch is a movable name pointing to a
commit; creating one does not duplicate the entire project. \code{HEAD} normally
refers to the branch currently checked out.

\begin{center}
\begin{tikzpicture}[
  commit/.style={circle,draw=RiceBlue,fill=RiceLightBlue,minimum size=7mm,font=\scriptsize},
  labelnode/.style={font=\small,align=left}
]
  \node[commit] (a) at (0,0) {A};
  \node[commit] (b) at (2,0) {B};
  \node[commit] (c) at (4,0) {C};
  \node[commit] (d) at (6,0.9) {D};
  \node[commit] (e) at (6,-0.9) {E};
  \draw[flowarrow] (a) -- (b);
  \draw[flowarrow] (b) -- (c);
  \draw[flowarrow] (c) -- (d);
  \draw[flowarrow] (c) -- (e);
  \node[labelnode,right=3mm of d] {\code{main}};
  \node[labelnode,right=3mm of e] {\code{feature}\\\code{HEAD}};
\end{tikzpicture}
\end{center}

In this graph, both branches share commits A through C. Work on \code{main}
produced D; work on \code{feature} produced E. A later merge can integrate the
histories. Switching branches moves \code{HEAD} and updates the working tree to
the selected snapshot. Git may refuse to switch when uncommitted changes would
be overwritten; that refusal protects evidence.

Commit identifiers are content-derived hexadecimal names, usually displayed in
abbreviated form. A branch name is convenient and movable; a commit identifier
names one historical state. Results that require durable reproducibility should
record a commit identifier or release tag, not only ``main,'' because
\code{main} will move.

\section{First-time setup}

Install Git using the current instructions for your operating system, then open
the repository folder in VS Code and use \textbf{Terminal: New Terminal}. The
integrated terminal may be PowerShell on Windows and a shell such as zsh or bash
on macOS or Linux. The Git commands in this appendix are the same in each; we
avoid shell-specific commands for creating or editing files.

Verify Git and set the identity written into future commits:

\begin{lstlisting}[language=bash,numbers=none]
git --version
git config --global user.name "Your Name"
git config --global user.email "you@example.edu"
git config --global init.defaultBranch main
git config --global --list
\end{lstlisting}

The name and email are commit metadata, not GitHub authentication. They should
identify your work accurately. If email privacy matters, use the no-reply address
provided by your GitHub account and follow the account's current email settings.
Repository-specific settings can override global settings by omitting
\code{--global} while inside that repository.

Authentication is a separate boundary. GitHub supports HTTPS and SSH remote
URLs. For HTTPS, a credential manager or \code{gh auth login} can complete a
browser-based sign-in; account passwords are not used as Git credentials. SSH
uses a private key on your computer and a corresponding public key registered
with GitHub. Follow GitHub's current setup instructions rather than pasting a
token into a command, file, notebook, or remote URL.

\begin{warningbox}
Never commit a password, API key, access token, private SSH key, database URL, or
credential file. Deleting it in a later commit does not remove it from earlier
history. Revoke or rotate an exposed credential immediately, then follow the
repository owner's incident and history-cleanup process.
\end{warningbox}

\section{Worked lesson 1: build a local history}

This lesson requires no GitHub account and makes no network changes. Create an
empty folder named \code{git-practice} outside the course repository, open that
folder in a separate VS Code window, and open its integrated terminal.

\subsection*{Step 1: initialize and inspect}

\begin{lstlisting}[language=bash,numbers=none]
git init -b main
git status
\end{lstlisting}

\code{git init} creates the hidden local repository. The \code{-b main} option
names the initial branch. No project snapshot exists until the first commit.

\subsection*{Step 2: create the first artifact}

In VS Code, create \code{README.md} containing:

\begin{lstlisting}[numbers=none,commentstyle=\color{CodeComment}]
# Probability experiment

This project records a small reproducible experiment.
\end{lstlisting}

Save the file, then inspect and stage it:

\begin{lstlisting}[language=bash,numbers=none]
git status
git diff -- README.md
git add README.md
git status
git diff --staged
\end{lstlisting}

Before \code{git add}, the file is untracked, so ordinary \code{git diff} may not
display its contents. After staging, \code{git diff --staged} shows the exact
snapshot proposed for the commit.

\subsection*{Step 3: make and inspect a commit}

\begin{lstlisting}[language=bash,numbers=none]
git commit -m "Document the probability experiment"
git status
git log --oneline --decorate
git show --stat --oneline HEAD
\end{lstlisting}

The commit message completes the sentence ``If applied, this commit will...''
with a concise action. \code{HEAD} refers to the current commit. A clean status
now means the working tree and staging area match that commit.

\subsection*{Step 4: separate two changes}

Create \code{experiment.py} in VS Code:

\begin{lstlisting}
outcomes = [True, False, True, True, False]
estimated_probability = sum(outcomes) / len(outcomes)
print(estimated_probability)
\end{lstlisting}

Also add a sentence to \code{README.md} explaining that the estimate is the
fraction of observed successes. Save both files and run:

\begin{lstlisting}[language=bash,numbers=none]
git status --short
git diff
git add experiment.py
git diff --staged
git commit -m "Add a Bernoulli probability estimate"
git status --short
\end{lstlisting}

Only \code{experiment.py} entered that commit. The README edit remains in the
working tree. Now review, stage, and commit it separately:

\begin{lstlisting}[language=bash,numbers=none]
git diff -- README.md
git add README.md
git diff --staged
git commit -m "Explain the probability estimate"
git log --oneline --decorate --graph --all
\end{lstlisting}

This is the purpose of staging: construct one reviewable snapshot from selected
work. \code{git add .} is convenient but broad; named paths make intent visible
and reduce accidental inclusion.

\begin{checkpointbox}
Edit \code{experiment.py}, stage it, and then edit it again. Use \code{git status},
\code{git diff}, and \code{git diff --staged} to identify which version is in
each state. Explain what a commit would record before running any commit command.
\end{checkpointbox}

\section{Design coherent commits}

A useful commit is a review and recovery unit. It should express one coherent
intent, include the tests and documentation required by that intent, and avoid
unrelated formatting or generated noise. ``One intent'' does not mean one file:
a feature may properly change source, tests, documentation, and configuration
together.

Before committing, ask:

\begin{itemize}
  \item Can I state the purpose in one specific sentence?
  \item Does every staged line contribute to that purpose?
  \item Did I include tests, documentation, schema, or migration changes needed
  for the contract?
  \item Did I inspect both unstaged and staged differences?
  \item Did I run the relevant checks from the intended environment?
  \item Is any secret, private data, large generated file, or local path present?
\end{itemize}

Prefer an imperative subject such as ``Validate probability bounds'' over
``changes,'' ``updates,'' or ``worked on notebook.'' The body of a larger commit
can explain why the change is needed, important constraints, and consequences
that are not obvious from the diff. A commit message should explain intent; it
should not repeat filenames already visible in the snapshot.

\section{Ignore local and generated state deliberately}

A \code{.gitignore} file names untracked paths that Git should normally omit
from status and staging. This repository ignores the project virtual
environment, Python caches, notebook checkpoints, LaTeX intermediates, editor
state, local secrets, and most local data. That policy keeps the repository
small and prevents machine-specific artifacts from obscuring meaningful change.

\begin{lstlisting}[numbers=none]
.venv/
__pycache__/
.ipynb_checkpoints/
.env
data/raw/*
data/processed/*
\end{lstlisting}

Ignoring is not security. It reduces accidental staging but does not protect a
file that was already tracked, explicitly forced into staging, copied elsewhere,
or exposed outside Git. Use \code{git check-ignore -v PATH} to determine which
rule ignores a path. Use \code{git ls-files PATH} to determine whether Git already
tracks it.

Git is optimized for source-like artifacts, not arbitrary data storage. Do not
commit virtual environments, dependency downloads, database files, model
checkpoints, or large raw datasets merely for convenience. Store reproducible
instructions, schemas, small teaching fixtures, checksums, and controlled
references instead. When a project genuinely requires versioned large binaries,
choose an approved artifact store or Git LFS and document the retrieval and
retention policy. GitHub currently warns for regular Git files above 50 MiB and
blocks files above 100 MiB; institutional and repository policies may be
stricter.

\section{Worked lesson 2: branches isolate change}

Return to \code{git-practice} with a clean working tree. Create and switch to a
branch in one command:

\begin{lstlisting}[language=bash,numbers=none]
git status
git switch -c feature/validate-probabilities
git branch --show-current
git log --oneline --decorate --graph --all
\end{lstlisting}

Edit \code{experiment.py} so that invalid values are rejected. One possible
version is:

\begin{lstlisting}
def mean_probability(values: list[float]) -> float:
    if not values:
        raise ValueError("values must not be empty")
    if any(value < 0 or value > 1 for value in values):
        raise ValueError("probabilities must lie in [0, 1]")
    return sum(values) / len(values)


print(mean_probability([1.0, 0.0, 1.0, 1.0, 0.0]))
\end{lstlisting}

Review and commit:

\begin{lstlisting}[language=bash,numbers=none]
git diff --check
git diff -- experiment.py
git add experiment.py
git diff --staged
git commit -m "Validate probability inputs"
git log --oneline --decorate --graph --all
\end{lstlisting}

Switch to the main branch and inspect the file:

\begin{lstlisting}[language=bash,numbers=none]
git switch main
git log --oneline --decorate --graph --all
git switch feature/validate-probabilities
\end{lstlisting}

The function disappears on \code{main} and returns on the feature branch because
the working tree follows \code{HEAD}. No work was lost. The branch names point
to different commits.

To integrate the completed branch locally:

\begin{lstlisting}[language=bash,numbers=none]
git switch main
git merge --ff-only feature/validate-probabilities
git log --oneline --decorate --graph --all
git branch -d feature/validate-probabilities
\end{lstlisting}

A fast-forward moves \code{main} to the descendant commit without creating a
merge commit. Deleting the merged branch removes a movable name, not the commits
now reachable from \code{main}. In a shared GitHub repository, integration
normally happens through a reviewed pull request instead of a local merge.

\section{Connect a repository to GitHub}

There are two common starting paths.

\textbf{Clone an existing repository.} On GitHub, copy the HTTPS or SSH URL from
the repository's Code menu. In a parent directory, run:

\begin{lstlisting}[language=bash,numbers=none]
git clone <repository-url>
cd <repository-directory>
git remote -v
git status
\end{lstlisting}

Cloning creates a new directory, copies reachable history, checks out a branch,
and records the source remote as \code{origin}. Angle-bracketed names above are
placeholders; replace them, including the brackets, with actual values.

\textbf{Publish an existing local repository.} First create an empty GitHub
repository without an automatically generated README or license, then connect
the local history:

\begin{lstlisting}[language=bash,numbers=none]
git remote add origin <repository-url>
git remote -v
git push -u origin main
\end{lstlisting}

The \code{-u} option records \code{origin/main} as the upstream for local
\code{main}, allowing later \code{git push} and status comparisons to infer the
destination. Pushing transfers commits and branch references; it does not upload
uncommitted files.

If you lack write permission to a repository, GitHub's fork model can create a
repository under your account from which you propose a pull request. A fork is
a GitHub repository relationship, not the same concept as a local branch.

\section{The daily branch and pull-request workflow}

GitHub flow is a lightweight branch-based workflow. For this course, use the
following conservative sequence unless an assignment specifies otherwise.

\subsection*{1. Begin from an inspected, current base}

\begin{lstlisting}[language=bash,numbers=none]
git status
git switch main
git fetch origin
git status -sb
git pull --ff-only
\end{lstlisting}

\code{git fetch} downloads remote commits and updates remote-tracking names such
as \code{origin/main}; it does not rewrite your working files. Inspect before
integrating. \code{git pull --ff-only} permits only an unambiguous fast-forward
and stops if local and remote history diverged.

\subsection*{2. Create one purpose-specific branch}

\begin{lstlisting}[language=bash,numbers=none]
git switch -c feature/add-calibration-check
\end{lstlisting}

Use a short descriptive name. Separate unrelated work into separate branches so
review, delay, or rejection of one change does not entangle another.

\subsection*{3. Work in small verified steps}

Edit in VS Code. Frequently inspect:

\begin{lstlisting}[language=bash,numbers=none]
git status --short
git diff
uv run pytest -q
uv run ruff check src tests scripts
\end{lstlisting}

Stage explicit paths and review the proposed commit:

\begin{lstlisting}[language=bash,numbers=none]
git add src/rice_dsm tests
git diff --staged --check
git diff --staged
git commit -m "Validate calibration ranges"
\end{lstlisting}

\subsection*{4. Publish the branch}

\begin{lstlisting}[language=bash,numbers=none]
git push -u origin feature/add-calibration-check
\end{lstlisting}

The feature branch now exists locally and on GitHub. Future commits pushed to
that branch automatically update its pull request.

\subsection*{5. Open a pull request}

On GitHub, choose the feature branch as the compare branch and \code{main} as the
base branch. The pull request asks to integrate the feature history into the
base; it is not merely a notification that coding is finished. A strong
description answers:

\begin{itemize}
  \item What problem or scientific need does this change address?
  \item What behavior, assumptions, data, or interfaces changed?
  \item What evidence was run, and what does it not prove?
  \item What should the reviewer inspect especially carefully?
  \item What could fail, and how would the team recover?
\end{itemize}

Use a draft pull request when the design or evidence needs early discussion.
Read the Files changed tab as the integrated proposal, not merely commit by
commit. Inspect automated checks and respond to review with additional commits;
do not resolve a conversation until the concern is actually addressed or moved
to an explicit follow-up issue.

\subsection*{6. Merge and clean up}

After required review and checks pass, use the repository's selected merge
strategy. GitHub may offer merge commits, squash merging, or rebasing; the team
should choose deliberately and consistently. Delete the merged remote branch,
then synchronize locally:

\begin{lstlisting}[language=bash,numbers=none]
git switch main
git pull --ff-only
git branch -d feature/add-calibration-check
\end{lstlisting}

\begin{professionalbox}
Never push directly to a protected default branch simply because you have
permission. Branch protection, review, and required checks make integration a
visible decision. Administrator ability is not evidence that bypassing the
process is appropriate.
\end{professionalbox}

\section{Review is a technical skill}

A reviewer is not a final syntax checker. Review challenges the change at
several levels.

\begin{center}
\small
\begin{tabularx}{\textwidth}{@{}
  >{\raggedright\arraybackslash}p{0.20\textwidth}
  >{\raggedright\arraybackslash}X
  >{\raggedright\arraybackslash}X@{}}
\toprule
\textbf{Lens} & \textbf{Reviewer question} & \textbf{Possible evidence} \\
\midrule
Intent & Does the diff solve the stated problem without unrelated change? &
Issue, pull-request description, coherent commits \\
Contract & Are inputs, outputs, errors, schemas, units, and compatibility clear? &
Types, docstrings, validation, interface tests \\
Correctness & Does behavior hold for ordinary, boundary, and failure cases? &
Focused tests, properties, independent calculation \\
Data and science & Are provenance, leakage, uncertainty, and interpretation
addressed? & Data checks, split design, evaluation, domain review \\
Operations & Can the change be observed, bounded, secured, and recovered? &
Logs, limits, permissions, migration and rollback plan \\
Maintainability & Can a future contributor understand and safely change it? &
Names, structure, documentation, small diff \\
\bottomrule
\end{tabularx}
\end{center}

Review comments should identify the observation and its consequence. Prefer
``This parser accepts a missing unit, so Celsius and Fahrenheit records can
become indistinguishable; could validation reject an absent unit?'' over ``This
looks wrong.'' Ask questions when context is missing. Mark nonblocking
suggestions as such. Authors should explain disagreement with evidence rather
than silently obeying or dismissing feedback.

\section{Synchronize without guessing}

Three commands are often confused:

\begin{description}
  \item[\code{git fetch origin}] Downloads remote objects and updates local
  remote-tracking references. It does not integrate them into the current branch.
  \item[\code{git merge origin/main}] Integrates the fetched \code{origin/main}
  history into the current branch, possibly by fast-forward or merge commit.
  \item[\code{git pull}] Fetches and then integrates according to configuration
  and options. Because it combines steps, inspect status and use an explicit
  policy such as \code{--ff-only} when that is the intended result.
\end{description}

Useful comparisons after fetching include:

\begin{lstlisting}[language=bash,numbers=none]
git status -sb
git log --oneline --decorate --graph --all -n 20
git log main..origin/main --oneline
git diff main...origin/main
\end{lstlisting}

The two-dot log above asks which commits are reachable from \code{origin/main}
but not local \code{main}. The three-dot diff compares each side from their
merge base and is useful for reviewing branch work. These notations are precise
but easy to reverse; state the question before trusting the output.

When a shared base moves while your feature branch is open, fetch first. The
repository may prefer merging \code{origin/main} into the feature or rebasing
the feature onto it. Rebasing rewrites commit identities and therefore requires
coordination; do not rebase commits that other people may have based work upon
unless the team workflow explicitly permits it.

\section{Worked lesson 3: create and resolve a conflict}

A conflict is not Git corruption. It means two histories made incompatible
changes and Git refuses to invent the intended result. Practice locally in
\code{git-practice} while its status is clean.

\subsection*{Step 1: create competing histories}

Create a branch from \code{main}:

\begin{lstlisting}[language=bash,numbers=none]
git switch main
git switch -c experiment/use-observed-rate
\end{lstlisting}

Edit one specific README sentence to say ``The reported estimate is the observed
success rate.'' Commit it:

\begin{lstlisting}[language=bash,numbers=none]
git add README.md
git commit -m "Define the estimate as an observed rate"
\end{lstlisting}

Return to \code{main}, create another branch, and change the same original
sentence to ``The reported estimate is a model probability.''

\begin{lstlisting}[language=bash,numbers=none]
git switch main
git switch -c experiment/use-model-probability
git add README.md
git commit -m "Define the estimate as a model probability"
\end{lstlisting}

\subsection*{Step 2: attempt integration}

While on \code{experiment/use-model-probability}, run:

\begin{lstlisting}[language=bash,numbers=none]
git merge experiment/use-observed-rate
git status
\end{lstlisting}

Git should report a content conflict. VS Code will show a merge editor or the
underlying markers:

\begin{lstlisting}[numbers=none]
  <<<<<<< HEAD
The reported estimate is a model probability.
  =======
The reported estimate is the observed success rate.
  >>>>>>> experiment/use-observed-rate
\end{lstlisting}

\code{HEAD} is the checked-out side; the other label names the merged side. The
final answer need not be either side verbatim. Decide the domain meaning first.
For this example, replace the entire marked region with a scientifically precise
sentence such as ``The script reports an observed success rate; it estimates a
probability only under stated sampling assumptions.'' Remove all markers.

\subsection*{Step 3: validate and complete}

\begin{lstlisting}[language=bash,numbers=none]
git diff --check
git diff
git add README.md
git status
git commit -m "Reconcile the probability interpretation"
git log --oneline --decorate --graph --all
\end{lstlisting}

Staging the file marks the conflict as resolved because you have selected its
final content. The merge commit records both parent histories. Run relevant
tests after resolution: choosing text that removes markers does not prove the
combined behavior is correct.

If you realize the merge began from the wrong state and have not committed its
resolution, \code{git merge --abort} normally returns to the pre-merge state.
Inspect \code{git status} before and after. Do not use destructive cleanup
commands merely to make the conflict display disappear.

\begin{checkpointbox}
Explain why Git could not safely choose between ``observed rate'' and ``model
probability.'' Then identify the scientific assumption introduced by the
resolved sentence. The conflict was textual, but the correct resolution required
domain reasoning.
\end{checkpointbox}

\section{Recover conservatively}

Recovery begins by classifying the state: untracked, unstaged, staged,
committed locally, pushed publicly, or merged. A command appropriate for one
state may destroy evidence in another.

\begin{center}
\small
\begin{tabularx}{\textwidth}{@{}
  >{\raggedright\arraybackslash}p{0.23\textwidth}
  >{\raggedright\arraybackslash}p{0.27\textwidth}
  >{\raggedright\arraybackslash}X@{}}
\toprule
\textbf{Intent} & \textbf{Conservative command} & \textbf{Meaning and caution} \\
\midrule
Inspect current state & \code{git status}; \code{git diff};
\code{git diff --staged} & Read before changing anything \\
Remove a path from staging & \code{git restore --staged PATH} & Keeps the
working-tree edit but removes it from the proposed commit \\
Discard an unstaged tracked edit & \code{git restore PATH} & Overwrites
uncommitted content; inspect and preserve needed work first \\
Correct the most recent unpushed commit & \code{git commit --amend} & Replaces
that commit with a new identity; avoid after sharing \\
Counteract a shared commit & \code{git revert COMMIT} & Creates a new commit
whose change reverses the selected commit, preserving public history \\
Abandon an unfinished merge & \code{git merge --abort} & Attempts to restore the
pre-merge state; check status and preserve unrelated work \\
Locate recently moved references & \code{git reflog} & Diagnostic record that
can help an experienced user recover otherwise unreachable commits \\
\bottomrule
\end{tabularx}
\end{center}

Commands such as \code{git reset --hard}, \code{git clean -fd}, and force push
can permanently discard uncommitted work or rewrite shared history. They are
not routine fixes. Stop, inspect exact targets, copy irreplaceable files outside
the repository, and seek review before using them. Never follow a destructive
command copied from an error search without understanding which state it changes.

\begin{warningbox}
Git does not record ordinary uncommitted edits. If a file was never committed or
stashed, Git may be unable to recover it after a destructive restore, clean, or
reset. Version control complements backups and editor recovery; it is not a
substitute for either.
\end{warningbox}

\section{Notebooks, data, and generated artifacts}

Jupyter notebooks are JSON documents containing source, metadata, execution
counts, and possibly outputs. A tiny logical edit can therefore create a noisy
diff, and hidden kernel state can make a recorded output disagree with visible
source. Before committing a teaching or analysis notebook:

\begin{enumerate}
  \item select the intended project kernel;
  \item restart and run all cells from top to bottom;
  \item inspect warnings, outputs, execution order, and data paths;
  \item remove accidental exploration, secrets, and machine-specific metadata;
  \item retain outputs only when repository policy says they teach or document
  something useful;
  \item inspect the Git diff and run the notebook regression tests; and
  \item keep reusable logic in package modules with focused tests.
\end{enumerate}

Data requires a separate versioning design. Git can identify small immutable
fixtures, schemas, metadata, and transformation code. It is often the wrong
storage system for sensitive records, frequently changing databases, or large
scientific arrays. Record a dataset identifier, checksum, provenance, schema,
access procedure, and transformation revision so that Git history can refer to
the evidence without embedding it improperly.

Generated artifacts should have an explicit policy. If a PDF, figure, or model
is tracked, document why the built output belongs in review and how to reproduce
it. Otherwise ignore it and let CI or a release system build it. Avoid commits in
which meaningful source changes are hidden beneath hundreds of regenerated
lines.

\section{GitHub Actions turns history into automated evidence}

GitHub Actions is an automation platform attached to repository events. A
\term{workflow} is a YAML file under \code{.github/workflows}. An event such as
a push, pull request, manual dispatch, or schedule triggers one or more jobs.
Each job runs on a runner and contains ordered steps. A step may run a shell
command or invoke a reusable action.

This repository's \code{course-ci.yml} workflow runs on pull requests, pushes
to \code{main}, and manual dispatch. Restricting ordinary push runs to
\code{main} avoids duplicating the pull-request matrix for a feature branch. It
grants read-only repository contents permission, cancels obsolete runs for the
same pull request or reference, and evaluates the course on Linux, macOS, and
Windows. Each runner installs the pinned tools, reproduces the locked
environment, prepares the notebook kernel, checks style, builds the package,
and runs all package, repository, and notebook tests. A stable final job named
\code{CI gate} requires every operating-system job to pass.

The core structure is:

\begin{lstlisting}[numbers=none]
name: Course CI

on:
  push:
    branches: [main]
  pull_request:
  workflow_dispatch:

permissions:
  contents: read

jobs:
  course-quality:
    strategy:
      matrix:
        os: [ubuntu-latest, macos-latest, windows-latest]
    runs-on: ${{ matrix.os }}
    steps:
      - uses: actions/checkout@3d3c42e5aac5ba805825da76410c181273ba90b1
      - run: uv sync --locked
      - run: uv run ruff check src tests scripts
      - run: uv run pytest -q
\end{lstlisting}

The excerpt omits setup steps for focus; the checked-in workflow is the
executable authority. Third-party actions are code executed with the workflow's
permissions. Pinning an action to a reviewed commit identifier reduces the risk
that a moving tag silently changes behavior. Keep workflow permissions narrow,
avoid exposing secrets to untrusted changes, set timeouts and concurrency
deliberately, and review dependency updates.

A green check has a precise meaning: the configured workflow passed for that
revision under those runner conditions. It does not prove the tests are
sufficient, the statistical design is valid, or deployment is safe. A red check
is evidence to inspect, not a button to rerun until it turns green. Open the job,
find the first meaningful failure, reproduce it locally when possible, repair
the cause, and push a new commit.

The repository also contains path-specific textbook, dependency-review, and
labeling workflows, plus two publication paths. Relevant merges to \code{main}
deploy the compiled textbook through a static course site. Matching annotated
\code{course-vX.Y.Z} tags build a checksummed release bundle and pause at a
protected approval environment before publication. Appendix~\ref{app:cicd}
traces those paths as a complete worked example and distinguishes integration,
delivery, deployment, and release.

\section{A command map by intention}

Memorize questions before commands.

\begin{longtable}{@{}
  >{\raggedright\arraybackslash}p{0.31\textwidth}
  >{\raggedright\arraybackslash}p{0.58\textwidth}@{}}
\toprule
\textbf{Question or intention} & \textbf{Command} \\
\midrule
\endhead
Where am I and what changed? & \code{git status}; \code{git status --short} \\
Which branch is checked out? & \code{git branch --show-current} \\
What is unstaged? & \code{git diff} \\
What is staged for the next commit? & \code{git diff --staged} \\
Does the diff contain whitespace errors? & \code{git diff --check};
\code{git diff --staged --check} \\
Select a path for the next commit & \code{git add PATH} \\
Unstage while keeping the working edit & \code{git restore --staged PATH} \\
Record the staged snapshot & \code{git commit -m "Imperative summary"} \\
Inspect compact history & \code{git log --oneline --decorate --graph --all} \\
Inspect one commit & \code{git show COMMIT} \\
Create and switch to a branch & \code{git switch -c BRANCH} \\
Switch to an existing branch & \code{git switch BRANCH} \\
List local and remote-tracking branches & \code{git branch -avv} \\
List configured remotes & \code{git remote -v} \\
Download remote history without integrating & \code{git fetch origin} \\
Fast-forward the current branch from its upstream & \code{git pull --ff-only} \\
Publish a new branch and record its upstream & \code{git push -u origin BRANCH} \\
Integrate only if a fast-forward is possible & \code{git merge --ff-only BRANCH} \\
Abort an unfinished merge & \code{git merge --abort} \\
Create a public counteracting commit & \code{git revert COMMIT} \\
Determine why a path is ignored & \code{git check-ignore -v PATH} \\
Determine whether a path is tracked & \code{git ls-files PATH} \\
\bottomrule
\end{longtable}

Run \code{git COMMAND --help} for the complete installed manual. Options and
defaults matter; read the command you are about to execute rather than assembling
fragments from unrelated examples.

\section{Principles to carry into professional work}

\begin{enumerate}
  \item \textbf{Inspect before changing state.} Status and diffs are part of the
  operation, not optional cleanup.
  \item \textbf{Commit one coherent intent.} Make history readable enough to
  review, test, and reverse selectively.
  \item \textbf{Branch by purpose.} Keep unrelated work and review conversations
  independent.
  \item \textbf{Never use the repository as a secret store.} Ignore files for
  convenience; use credential systems for security.
  \item \textbf{Keep data identity separate from data bulk.} Version schemas,
  provenance, code, and identifiers; store sensitive or large evidence in an
  appropriate controlled system.
  \item \textbf{Treat pull requests as reasoning records.} Explain intent,
  evidence, risk, and review decisions.
  \item \textbf{Protect shared history.} Prefer additive correction with revert;
  coordinate any history rewrite.
  \item \textbf{Interpret CI precisely.} A check supports only the contracts it
  actually exercises.
  \item \textbf{Make recovery proportional and conservative.} Classify the state,
  preserve evidence, then choose the narrowest operation.
  \item \textbf{Name released artifacts by immutable revisions.} Moving branch
  names are convenient for development, not sufficient provenance for a result.
\end{enumerate}

\section{Exercises}

\subsection*{Foundational}

\begin{enumerate}
  \item Draw the working tree, staging area, local repository, and remote
  repository. Place \code{add}, \code{commit}, \code{push}, and \code{fetch} on
  the correct transitions. Explain what each command does not do.
  \item Distinguish a branch, a remote-tracking branch, a fork, and a pull
  request. Give a situation in which all four appear.
  \item A file is listed as both staged and modified. Explain how this state can
  arise and which two diff commands reveal its two versions.
\end{enumerate}

\subsection*{Applied}

Complete Worked Lessons 1 through 3 in a disposable practice repository. Submit
the output of \code{git log --oneline --decorate --graph --all}, one staged diff,
and a short explanation of the conflict resolution. Do not submit the hidden
\code{.git} directory.

Inspect this course repository's \code{.gitignore} and
\code{.github/workflows/course-ci.yml}. For five ignored path groups, explain
the failure prevented. For every CI step, name the contract it checks and one
important failure it cannot detect.

\subsection*{Design}

Design a GitHub workflow for a four-person scientific software team. Specify
branch naming, commit expectations, required pull-request evidence, reviewers,
CI checks, data and secret policy, merge strategy, release tags, and emergency
recovery authority. Include one case where a green CI run should still block
merge and one case where rewriting history is justified but requires
coordination.

\begin{chapterreview}
\textbf{Key ideas.} Git records selected project snapshots and their history;
GitHub organizes shared review and automation around that history. The staging
area constructs the next commit. Branches are movable commit references, not
copied folders. Pull requests expose proposed integration, discussion, and CI
evidence. Recovery begins by identifying state and choosing the narrowest safe
operation.

\textbf{Vocabulary.} repository; working tree; staging area; commit; branch;
HEAD; remote; remote-tracking branch; clone; fetch; push; merge; conflict; fork;
pull request; workflow; runner; revert.

\textbf{Explain aloud.} Why are ``my file is saved,'' ``my change is committed,''
and ``my branch is on GitHub'' three different claims?
\end{chapterreview}

\section{Official references}

Interfaces and recommendations evolve. Verify account, authentication, and
workflow details against the current official documentation:

\begin{itemize}
  \item \href{https://git-scm.com/book/en/v2/Git-Basics-Recording-Changes-to-the-Repository}{Pro Git: Recording Changes to the Repository}
  \item \href{https://git-scm.com/book/en/v2/Git-Branching-Branches-in-a-Nutshell}{Pro Git: Branches in a Nutshell}
  \item \href{https://git-scm.com/docs}{Git command reference}
  \item \href{https://docs.github.com/en/get-started/using-github/github-flow}{GitHub flow}
  \item \href{https://docs.github.com/en/authentication}{GitHub authentication}
  \item \href{https://docs.github.com/en/pull-requests}{Pull requests and review}
  \item \href{https://docs.github.com/en/actions}{GitHub Actions}
  \item \href{https://docs.github.com/en/repositories/working-with-files/managing-large-files}{GitHub large-file guidance}
\end{itemize}

\conceptcheck{A teammate says, ``Everything is safe because I pushed it to a
branch and CI passed.'' Which claims are supported, which are not, and what
additional evidence would you request before merging?}
